\subsection{Phạm vi MVP}
\label{sec:mvp-scope}

MVP của OneSpace gồm toàn bộ tính năng trong bảng~\ref{tab:mvp-scope-full}. Bốn phân hệ đều phải có mặt. Cả năm trụ Winning Zone W1--W5 phải chạy được end-to-end. Mã tính năng tham chiếu bảng~\ref{tab:feature-pool}; mỗi tính năng đều truy vết được về pain point ở mục~\ref{sec:problem-statement} hoặc nhân tố Winning Zone ở bảng~\ref{tab:winning-zone}.

\begin{table}[H]
\centering
\caption{Phạm vi MVP của OneSpace}
\label{tab:mvp-scope-full}
\begin{tabularx}{\textwidth}{>{\raggedright\arraybackslash}p{1.8cm} >{\raggedright\arraybackslash}X}
    \toprule
    \textbf{Phân hệ} & \textbf{Tính năng trong MVP} \\
    \midrule
    Video & Xem video (F-VID-01); feed theo sở thích (F-VID-02); tìm kiếm (F-VID-03); PIP (F-VID-08); tua nhanh double-tap (F-VID-14); watchlist + tiếp tục từ chỗ dừng (F-VID-04, F-VID-05); chọn chất lượng Auto/1080p/720p/360p (F-VID-06); audio-only (F-VID-09); tải video tùy bản quyền (F-VID-07); trích mp3 (F-VID-10); chia sẻ link (F-VID-11); màn hồ sơ nghệ sĩ (F-VID-12); playlist video (F-VID-13); thẻ quiz cuối video (F-VID-15) \\
    Music & Player đầy đủ (F-MUS-01); dock mini player (F-MUS-02); nhạc nền ngoài app (F-MUS-03); playlist + queue (F-MUS-04, F-MUS-05); tìm kiếm nhạc (F-MUS-06); kho nhạc cá nhân (F-MUS-07); gợi ý theo sở thích (F-MUS-08); siêu playlist (F-MUS-09); audio focus (F-MUS-10); tải nhạc tùy bản quyền (F-MUS-12) \\
    News & Feed theo sở thích (F-NEWS-01); trình đọc + đổi cỡ chữ (F-NEWS-02); tìm kiếm (F-NEWS-04); nghe tin TTS (F-NEWS-05); sắp xếp mới nhất/phổ biến (F-NEWS-06); nghe nhạc khi đọc (F-NEWS-07); quiz từ bài báo (F-NEWS-08); gợi ý bài liên quan (F-NEWS-09); thẻ nổi khi rời tab (F-NEWS-10) \\
    Games & Hub minigame (F-GAME-01); đoán bài hát (F-GAME-02); quiz tin tức (F-GAME-03); quiz kiến thức kiểu Ai là triệu phú (F-GAME-04); Snake (F-GAME-05); Beat Tiles --- nốt rơi theo BPM thật của bài đang phát (F-GAME-06); bảng xếp hạng (F-GAME-07); game + video hai nửa màn (F-GAME-08); điểm thưởng khi thắng (F-GAME-09); 2048 (F-GAME-10); Tic-Tac-Toe (F-GAME-11); thử thách hằng ngày (F-GAME-12); âm thanh và hiệu ứng game (F-GAME-13) \\
    Dùng chung & Tài khoản (F-COM-01); onboarding chọn sở thích (F-COM-02); hồ sơ + đổi avatar + khung (F-COM-03, F-COM-09); hệ điểm XP (F-COM-04); nhiệm vụ hằng ngày (F-COM-05); streak (F-COM-06); pet (F-COM-07); huy hiệu (F-COM-08); wrapped (F-COM-10); tìm kiếm xuyên nội dung (F-COM-11); nhắc nhiệm vụ (F-COM-12); cài đặt (F-COM-13); popup xung đột âm thanh (F-COM-14); cửa hàng đổi quà (F-COM-15); offline (F-COM-16); chế độ khách (F-COM-18) \\
    \bottomrule
\end{tabularx}
\end{table}

MVP bao gồm đủ sáu trụ Winning Zone W1--W6: trụ W6 (wrapped~--~tổng kết tháng, F-COM-10) nằm trong MVP và đã có màn tổng kết chi tiết từng mảng kèm chia sẻ.

\textbf{Giai đoạn v2 (sau khi MVP hoàn tất).} Các tính năng còn lại chia theo lý do:

\begin{itemize}
    \item Equalizer (F-MUS-11): giá trị thấp so với công sức.
    \item Thuật toán gợi ý học máy thật: MVP chỉ lọc theo sở thích khai báo.
    \item Thanh toán, gói Premium và quảng cáo: v1 phát hành miễn phí để đo hành vi người dùng.
\end{itemize}

\textbf{Chỉ số kiểm chứng MVP.} MVP phải trả lời câu hỏi: người dùng có sẵn sàng rời các app chuyên biệt để dùng OneSpace không? Bảng~\ref{tab:mvp-metrics} liệt kê các chỉ số đo lường.

\begin{longtable}{>{\raggedright\arraybackslash}p{4.2cm} >{\raggedright\arraybackslash}p{8cm} >{\raggedright\arraybackslash}p{3.2cm}}
\caption{Chỉ số kiểm chứng MVP}
\label{tab:mvp-metrics}\\
\toprule
\textbf{Chỉ số} & \textbf{Trả lời câu hỏi gì} & \textbf{Ngưỡng đề xuất} \\
\midrule
\endfirsthead
\toprule
\textbf{Chỉ số} & \textbf{Trả lời câu hỏi gì} & \textbf{Ngưỡng đề xuất} \\
\midrule
\endhead
\bottomrule
\endfoot
Retention ngày 7 & Giá trị sản phẩm có giữ chân người dùng không? & $\geq$ 4\% theo \texttt{BIZ-01} (mục~\ref{sec:quality-pillars}) \\
\midrule
Tỉ lệ phiên dùng từ 2 phân hệ trở lên & Người dùng có chuyển tab không? & $\geq$ 40\% số phiên \\
\midrule
Số lượt chuyển tab không qua màn hình chính & Chuyển tab có mượt không (dock, PIP)? & Theo dõi, chưa cần ngưỡng \\
\midrule
Adoption tính năng USP trong tuần đầu & Tính năng nào hấp dẫn nhất? & $\geq$ 30\% người dùng mới dùng từng tính năng \\
\midrule
Tỉ lệ giữ streak $\geq$ 7 ngày & Gamification có tạo thói quen không? & $\geq$ 15\% người dùng hoạt động \\
\midrule
Thời gian phiên trung bình & Mức độ sử dụng tổng thể & $\geq$ 15 phút/người/ngày \\
\midrule
Crash-free và thời gian mở app & Chất lượng kỹ thuật & Crash-free $\geq$ 99,5\%, mở app $\leq$ 2 giây (\texttt{UX-03}, \texttt{UX-07}) \\
\end{longtable}

\textbf{Hành vi khi chuyển giữa các phân hệ.} Bảng~\ref{tab:switch-matrix} định nghĩa điều gì xảy ra khi người dùng chuyển tab, bằng vuốt ngang hoặc chạm thanh tab. Ma trận phủ \textbf{đủ 12 chiều chuyển} giữa bốn phân hệ A--D (mỗi phân hệ sang ba phân hệ còn lại, hai chiều \textit{không đối xứng}: A$\rightarrow$B và B$\rightarrow$A hành xử khác nhau), cộng bốn case đặc biệt C13--C16. Cột cuối ghi quyết định ảnh hưởng đến khối lượng code của Phần~\ref{sec:ui-ux-design} và Phần~\ref{sec:technical-architecture}. Ký hiệu: A = Video, B = Music, C = News, D = Games.

\begin{longtable}{>{\raggedright\arraybackslash}p{1.5cm} >{\raggedright\arraybackslash}p{8.6cm} >{\raggedright\arraybackslash}p{2.5cm} >{\raggedright\arraybackslash}p{2.6cm}}
\caption{Ma trận hành vi chuyển phân hệ}
\label{tab:switch-matrix}\\
\toprule
\textbf{Case} & \textbf{Người dùng thấy gì} & \textbf{Tính năng liên quan} & \textbf{Quyết định} \\
\midrule
\endfirsthead
\toprule
\textbf{Case} & \textbf{Người dùng thấy gì} & \textbf{Tính năng liên quan} & \textbf{Quyết định} \\
\midrule
\endhead
\bottomrule
\endfoot
C1 \newline A$\rightarrow$B \newline đang xem video, mở nhạc &
(1) Rời tab Video: video thu thành PIP góc dưới phải, vẫn còn tiếng. Sang tab Nhạc và chạm một bài: popup xuất hiện. \newline
(2) Popup hai lựa chọn \textit{Tắt tiếng video / Dừng video} kèm ô tick \textit{Không hỏi lại}; chọn áp dụng ngay cho lần đó — nhạc vẫn phát ở cả hai. \newline
(3) Không tick: lần xung đột sau vẫn hỏi lại; tick: ghi nhớ đúng lựa chọn vừa chọn, đặt lại được trong Cài đặt. \newline
(4) Dock hiện bài nhạc; PIP video và dock cùng tồn tại. &
F-VID-08, F-COM-14, F-MUS-02 &
Hỏi mỗi lần; tick Không hỏi lại để ghi nhớ \\
\midrule
C2 \newline B$\rightarrow$A \newline đang nghe nhạc, mở video &
(1) Mở video khi nhạc đang phát: popup hai lựa chọn --- \textit{Tắt tiếng video} hoặc \textit{Dừng nhạc} --- kèm ô tick \textit{Không hỏi lại}; tick để ghi nhớ. \newline
(2) Chọn tắt tiếng video: video mở có badge ``Tắt tiếng'', nhạc vẫn phát ở dock. \newline
(3) Chọn dừng nhạc: video phát tiếng, nhạc dừng ở dock, bấm Phát để nghe lại. &
F-VID-01, F-VID-09, F-MUS-02 &
Luồng mới phát chiếm quyền; nhạc nhường thầm lặng \\
\midrule
C3 \newline A$\rightarrow$C \newline đang xem video, đọc tin &
(1) Rời tab Video: video thành PIP góc, giữ tiếng theo chính sách đã chọn ở C1. \newline
(2) Đọc bài bình thường; cửa sổ nổi có nhãn ``Đang phát'' kèm ba nút Mở rộng, Tạm dừng, Đóng. \newline
(3) Chạm Mở rộng: quay lại trình phát đúng vị trí, lệch tối đa 2 giây. \newline
(4) Rời tab News: bài đang đọc thu thành thẻ ``Đang đọc'' nổi. &
F-VID-08, F-NEWS-10 &
PIP tự thu nhỏ khi rời tab; thẻ đọc nổi khi rời News \\
\midrule
C4 \newline C$\rightarrow$A \newline đang đọc tin, mở video &
(1) Đang đọc bài, mở video từ tab Video: video phát; thẻ ``Đang đọc'' nổi ở mọi tab (trong màn chơi thẻ tự ẩn để không che bàn chơi). \newline
(2) Chạm bất kỳ đâu trên thẻ: quay lại đúng chỗ đang đọc. &
F-NEWS-10, F-VID-01 &
Bài đọc dở không bao giờ mất \\
\midrule
C5 \newline A$\rightarrow$D \newline đang xem video, vào game &
(1) Vào game khi video đang phát: video chiếm khoảng hai phần năm phía trên kèm hàng điều khiển (phát/dừng, thanh tua, thời gian), game chiếm phần còn lại và nhận đủ cảm ứng. \newline
(2) Chia màn không chỉ có ở Games: cửa sổ nổi video có nút chia đôi, ghim video lên đỉnh màn hình và đẩy nội dung xuống, dùng được ở mọi tab. \newline
(3) Thoát game: video chạy tiếp; thoát video: game vẫn chạy. &
F-GAME-08, F-VID-08 &
Chia màn thay vì popup hỏi \\
\midrule
C6 \newline D$\rightarrow$A \newline đang trong game, mở video &
(1) Rời game: ván tự tạm dừng, trạng thái và điểm số được giữ. \newline
(2) Video phát; quay lại Games: ván dở tự hiện lại đúng trạng thái. &
F-GAME (pause/restore) &
Rời game luôn tạm dừng + lưu \\
\midrule
C7 \newline B$\rightarrow$C \newline đang nghe nhạc, đọc tin &
(1) Rời tab Music: nhạc vẫn phát; dock hiện ở feed và cả trong trình đọc bài. \newline
(2) Chạm dock mở trình phát đầy đủ, quay lại đọc tiếp bằng một chạm. \newline
(3) Chơi quiz cuối bài: nhạc vẫn phát. &
F-MUS-02, F-NEWS-07, F-NEWS-08 &
Dock không biến mất trong trình đọc \\
\midrule
C8 \newline C$\rightarrow$B \newline đang đọc tin, mở nhạc &
(1) Đang đọc, mở bài nhạc từ tab Music: nhạc phát; thẻ ``Đang đọc'' nổi theo dõi. \newline
(2) Quay lại bài bằng thẻ. &
F-MUS-01, F-NEWS-10 &
Đọc dở vẫn kiểm soát được \\
\midrule
C9 \newline B$\rightarrow$D \newline đang nghe nhạc, vào game &
(1) Nhạc vẫn phát trong game; âm hiệu ứng game tự giảm còn khoảng 25\% để không đè nguồn tiếng chính (mức âm chỉnh được trong Cài đặt). \newline
(2) Thắng ván: điểm cộng ngay, nhạc không ngắt. &
F-MUS-02, F-GAME-09, F-GAME-13 &
Âm game tự giảm khi nhạc phát \\
\midrule
C10 \newline D$\rightarrow$B \newline đang trong game, mở nhạc &
(1) Rời game: ván tạm dừng, trạng thái và điểm số được giữ. \newline
(2) Nhạc phát; quay lại game, ván dở tự hiện lại. &
F-MUS-01, F-GAME (pause/restore) &
Rời game luôn tạm dừng + lưu \\
\midrule
C11 \newline C$\rightarrow$D \newline đang đọc tin, vào game &
(1) Thẻ ``Đang đọc'' nổi (trong màn chơi thẻ tự ẩn để không che bàn chơi); game chạy bình thường. \newline
(2) Quay lại bài qua thẻ. &
F-NEWS-10, F-GAME-01 &
Thẻ đọc nổi dùng chung mọi luồng \\
\midrule
C12 \newline D$\rightarrow$C \newline đang trong game, đọc tin &
(1) Game tạm dừng; bài đọc mở; rời News thì thẻ ``Đang đọc'' nổi. &
F-GAME (pause/restore), F-NEWS-10 &
Rời game luôn tạm dừng + lưu \\
\midrule
C13 \newline B$\rightarrow$TTS \newline đang nghe nhạc, bấm \textit{Nghe} bài tin &
(1) Bấm nút Nghe trên bài báo: nếu nhạc đang phát, nhạc tự tạm dừng nhường. \newline
(2) TTS đọc theo đoạn kèm thanh tiến trình và nút dừng; nghe xong bật lại nhạc bằng dock. &
F-NEWS-05, F-MUS-10 &
TTS là luồng nói chính \\
\midrule
C14 \newline trích MP3 \newline (A$\rightarrow$B) &
(1) Nhấn giữ video $\rightarrow$ xác nhận \textit{Trích xuất} $\rightarrow$ tạo file âm thanh thật (TTS \texttt{synthesizeToFile} từng đoạn rồi ghép WAV) $\rightarrow$ bản ghi vào mục ``Từ video của bạn'' của kho Nhạc; phát như một bài nhạc thật, +5 XP lần đầu mỗi video. \newline
(2) Nhấn giữ bài báo: tương tự --- TTS tổng hợp thành file thật rồi ghép các đoạn lại. &
F-VID-10, F-MUS-07 &
Trích xuất dùng chung engine nhạc \\
\midrule
C15 \newline siêu playlist &
(1) Playlist là meta-playlist: chứa cả video lẫn bài hát, đếm kết hợp kiểu ``2 video $\cdot$ 3 bài hát''. \newline
(2) Phát tuần tự đan xen --- hết video tự chuyển sang bài hát và ngược lại; tiêu đề màn ``Danh sách phát: X''; nút Quay lại trở về danh sách playlist. &
F-MUS-09 &
Trong MVP (v1) \\
\midrule
C16 \newline ba luồng cùng lúc &
(1) Không hỗ trợ ở v1: giữ tối đa một nguồn tiếng chính và một PIP video. \newline
(2) Lý do: hiệu năng thiết bị phổ thông; ghi nhận cho v2. &
F-MUS-10, F-COM-14 &
Tối đa 1 nguồn tiếng + 1 PIP ở v1 \\
\end{longtable}

\textbf{Quy tắc xử lý xung đột âm thanh (áp dụng cho mọi case).}

\begin{itemize}
    \item Cuộc gọi và báo thức luôn cao nhất (audio focus của Android); mọi luồng của ứng dụng nhường ngay.
    \item Một nguồn tiếng chính tại một thời điểm: TTS làm nhạc tạm dừng (C13), âm hiệu ứng game tự giảm còn khoảng 25\% khi nhạc hoặc video đang chạy (C9).
    \item Mở nhạc khi video đang phát: hỏi (C1); mở video khi nhạc đang phát: hỏi (C2) --- cả hai chiều đều hỏi.
    \item Nhạc dock là luồng nền: hiện trên mọi tab, nhường ngay khi luồng chính cần tiếng, tiếp tục bằng một chạm.
    \item Mọi luồng đều thấy và điều khiển được từ dock mini player trong ứng dụng.
\end{itemize}

\textbf{Gián đoạn từ bên ngoài.}

\begin{itemize}
    \item Cuộc gọi đến: mọi luồng tạm dừng trong tối đa 0,5 giây; sau cuộc gọi nhạc tự phát lại, video hiện nút ``Tiếp tục'', game mở lại đúng ván dở.
    \item Rút tai nghe: nhạc và TTS dừng ngay (chuẩn Android), không phát loa ngoài; PIP video dừng kèm thông báo.
    \item Mất mạng khi đang xem video: khung hình đứng + tự thử lại sau 5--15 giây; tin đã tải đọc tiếp; game offline chơi tiếp, điểm lưu cục bộ rồi đồng bộ khi có mạng.
    \item Hệ điều hành đóng ứng dụng: mở lại khôi phục đủ 4 vị trí (bài nhạc, tin đang đọc, ván dở, video dở) trong 2 giây.
    \item Người dùng thoát ứng dụng: nhạc tiếp tục nếu đã bật nhạc nền (F-MUS-03); video, TTS dừng; game tạm dừng và giữ ván dở.
\end{itemize}

\textbf{Luồng lần đầu mở app.} Nếu người dùng mới không dùng được tính năng chính trong lần đầu thì mọi thiết kế phía sau vô nghĩa. Kịch bản:

\begin{enumerate}
    \item Mở app trong tối đa 2 giây (không bắt đăng nhập).
    \item Chọn tối thiểu 3 chủ đề (cho phép Bỏ qua).
    \item Vào Home thấy ngay ``Đề xuất hôm nay'': một video, một bài hát, một tin.
    \item Chạm một bài hát: dock mini hiện ra ở đáy mọi màn hình.
    \item Banner nhiệm vụ ngày đầu: \textit{Nghe 1 bài -- Đọc 1 tin -- Thắng 1 ván}.
    \item Mở bài tin đầu tiên: dock nhạc vẫn hiện trong trình đọc.
    \item Hoàn thành 3 nhiệm vụ: nhận trứng thú cưng, nở trong 2 giây.
    \item Cuối phiên: streak = 1; huy hiệu ``Chuỗi 7 ngày'\,hiện trong lưới huy hiệu.
\end{enumerate}

Chỉ số: \textbf{tỉ lệ người mới dùng ít nhất 3 trong 5 tính năng chính trong lần đầu} $\geq$ 60\%.

\textbf{Hệ thống điểm.} Bảng~\ref{tab:point-economy} xác định cách kiếm điểm và tiêu điểm.

\begin{longtable}{>{\raggedright\arraybackslash}p{7.6cm} >{\raggedright\arraybackslash}p{6.4cm}}
\caption{Cách kiếm điểm và tiêu điểm}
\label{tab:point-economy}\\
\toprule
\textbf{Kiếm điểm} & \textbf{Tiêu điểm} \\
\midrule
\endfirsthead
\toprule
\textbf{Kiếm điểm} & \textbf{Tiêu điểm} \\
\midrule
\endhead
\bottomrule
\endfoot
Thắng minigame theo từng game (xem dưới) & Cho thú cưng ăn: $-$30 (một bữa mỗi ngày) \\
Nhiệm vụ hằng ngày: +20 mỗi việc & Mở khoá game: $-$300 đến $-$500 \\
Quiz cuối video: 20 + 5$\times$(combo$-$1), tối đa +50 & Khung avatar: $-$200 đến $-$800 \\
Quiz từ bài báo: +15 mỗi câu đúng & Phụ kiện pet: $-$50 đến $-$100 \\
Quiz kiến thức: thang 5--60 điểm mỗi ván & (v2: đổi quà vật lý, voucher) \\
2048: điểm chia 20, chạm ô 2048 +100 &  \\
Thắng Tic-Tac-Toe: Dễ +10 / Vừa +20 / Khó +40 &  \\
Rắn săn mồi: +2 mỗi mồi &  \\
Beat Tiles: điểm chia 2 &  \\
Đoán bài hát: +10 mỗi câu đúng &  \\
Trích MP3: +5 lần đầu mỗi video &  \\
Thử thách hằng ngày: 50 + 5$\times$chuỗi ngày (tối đa 85) &  \\
\end{longtable}

Người chơi trung bình kiếm 110--150 điểm/ngày. Không bán điểm bằng tiền ở v1. Toàn bộ hệ số là hằng số cấu hình, chỉnh không cần build lại.

\textbf{Thú cưng ảo.} Thú cưng tiến hoá qua bốn giai đoạn theo tổng điểm tích luỹ: \textbf{Trứng} (0đ, nở sau khi hoàn thành 3 nhiệm vụ) $\rightarrow$ \textbf{Bé cưng} (500đ, gà con vàng) $\rightarrow$ \textbf{Thiếu niên} (2.000đ) $\rightarrow$ \textbf{Trưởng thành} (5.000đ, vẹt).

\textbf{Rủi ro và cách giảm thiểu.}

\begin{longtable}{>{\raggedright\arraybackslash}p{3.9cm} >{\centering\arraybackslash}p{1.1cm} >{\raggedright\arraybackslash}p{5.4cm} >{\raggedright\arraybackslash}p{4.2cm}}
\caption{Rủi ro và giảm thiểu}
\label{tab:risk-register}\\
\toprule
\textbf{Rủi ro} & \textbf{Mức} & \textbf{Nếu xảy ra} & \textbf{Giảm thiểu} \\
\midrule
\endfirsthead
\toprule
\textbf{Rủi ro} & \textbf{Mức} & \textbf{Nếu xảy ra} & \textbf{Giảm thiểu} \\
\midrule
\endhead
\bottomrule
\endfoot
Bản quyền trích mp3 & Cao & Bị khiếu nại, gỡ app & Chỉ trích từ catalog tự dựng; ghi rõ phạm vi nội dung tự dựng \\
TTS tiếng Việt kém & Trung bình & Trụ W3 nghe tệ & Chọn engine giọng Việt ổn; thử 20 bài; fallback hiện bản chữ \\
OEM tắt ứng dụng nền & Cao & Nhạc nền chết trên máy Xiaomi/Oppo & Foreground service đúng chuẩn; hướng dẫn người dùng tắt tối ưu pin; test trên 2--3 máy OEM \\
Game + video trên máy yếu & Trung bình & C5 giật, người bỏ ván & PIP tự hạ độ phân giải; test máy RAM 3 GB \\
Phạm vi phình to & Cao & MVP trễ, chất lượng giảm & Freeze danh sách MVP; mọi đề nghị mới vào danh sách v2 \\
Gian lận điểm & Trung bình & BXH vô nghĩa & v1 ghi điểm cục bộ; chống gian lận sau \\
Thiết bị Android đa dạng & Trung bình & Lỗi chỉ trên máy của giảng viên & Android 8.0+; test trên 3 máy cố định \\
Thành viên nghỉ giữa chừng & Trung bình & Module trễ & Mỗi module có người phụ trách + tài liệu; review chéo \\
\end{longtable}

\textbf{Quyền hệ thống.} Bảng~\ref{tab:permissions} liệt kê toàn bộ quyền ứng dụng cần.

\begin{table}[H]
\caption{Quyền hệ thống}
\label{tab:permissions}
\centering
\begin{tabularx}{\textwidth}{>{\raggedright\arraybackslash}p{3.4cm} >{\raggedright\arraybackslash}p{4.6cm} >{\raggedright\arraybackslash}p{3.7cm} >{\raggedright\arraybackslash}p{3.0cm}}
\toprule
\textbf{Quyền} & \textbf{Tính năng nào cần} & \textbf{Xin khi nào} & \textbf{Nếu từ chối} \\
\midrule
INTERNET & Mọi tính năng tải nội dung & Khai báo mặc định, không hỏi & Vào chế độ ngoại tuyến \\
POST\_\newline NOTIFICATIONS & Nhắc nhiệm vụ, điều khiển nhạc & Chỉ khi người dùng bật nhắc & Loại nhắc đó tắt; phần còn lại hoạt động bình thường \\
FOREGROUND\_\newline SERVICE & Nhạc nền (F-MUS-03) & Lần đầu rời app khi đang phát nhạc & Nhạc dừng khi rời app \\
\bottomrule
\end{tabularx}
\end{table}

\textbf{Thiết kế để mở rộng sau này.} Bốn quy tắc để MVP không chặn hướng phát triển của v2:

\begin{itemize}
    \item Không chèn quảng cáo vào giữa các lần chuyển tab. Sau này có quảng cáo thì đặt ở điểm dừng tự nhiên (kết thúc ván, màn nhiệm vụ).
    \item Điểm chỉ đổi hàng phi chức năng (pet, khung, mở game): không khóa nội dung. Sau này thêm Premium không ảnh hưởng người dùng miễn phí.
    \item Converter (TTS, trích mp3) là module riêng, engine thay được: đổi nhà cung cấp TTS không đụng UI.
    \item Nội dung quy về ``content ID'' chung: v2 thêm nguồn mới không phải sửa danh mục tính năng.
\end{itemize}

\textbf{Ước lượng khối lượng.} Media engine + dock + audio focus: 3--4 tuần-người (làm trước). Bốn phân hệ UI: 1,5--2 tuần-người mỗi phân hệ. Hệ điểm + nhiệm vụ + pet: 1,5--2. TTS: 1--2. Mỗi minigame: 1--2. PIP + ma trận case: 1--1,5. Tổng cho nhóm 3 người làm 8--10 tuần: vừa đủ nếu không phình phạm vi.
